% ===================================================================
%  TAULUKKO N:o 8 — Elonkorjuu. — Metsästys, viininkorjuu, kalastus.
%  Traduction 2026 d'apres texto/io, controlee sur texto/fr et
%  texto/sv. Consigne : voir texto/fi/10-taulukko-01.tex.
%
%  L'appel de note est dans le TITRE de la scene. LE RENVOI (13) DE
%  LA FAUX MANQUE A L'IDO ; la colonne le rend : ecart declare.
% ===================================================================

%%K t08-noto-1 noto
\VUnotes{143.2mm}{%
\textit{(*) Koska tuo sana ei ole täsmällinen : onko kyseessä pellavan korjuu,
viininkorjuu, omenien korjuu? Ehkä olisi hyvä käyttää sanaa «} \VUgras{mesono}
\textit{», jota ehdotettiin Mondo XIV, s. 242. Mutta tähän asti Akatemia ei ole
hyväksynyt uutta juurta, ja se odottaa keskusteluja tavanomaisen idistisen
ohjesäännön mukaan.}%
}

%%K t08-apar-1 apar
\VUsaut{5.0mm}
\VUtitre{15.0pt}{60}{TAULUKKO N:o 8}
\VUsaut{3.0mm}
\VUcentre{13.2pt}{90}{\textit{Ensimmäinen kohtaus.}}
\VUsaut{3.0mm}
\VUpk{10.2pt}{40}{Elonkorjuu (*).}
\VUsaut{4.0mm}
\VUpk{10.2pt}{40}{Maaseudun näkymiä.}
\VUsaut{3.0mm}

%%K t08-c1-01-1 p
1. --- \VUgras{Kesän} myötä maaseudun ulkonäkö on jälleen kerran muuttunut.
Siemen, joka uskottiin vaolle, on itänyt; vihreä taimi nousi maasta ja meni
tähkälle. Jyvä muodostui, sitten se kypsyi, ja nyt tarvitsee vain korjata sato.

%%K t08-c1-02-1 p
2. --- \VUgras{Niittykukat} ovat puhjenneet. \VUgras{Ruiskaunokki}
\textsuperscript{(1)}, \VUgras{unikko} \textsuperscript{(2)} ja
\VUgras{sormustinkukka} \textsuperscript{(3)} sekoittavat vehnäpeltojen
yksitoikkoiseen värisävyyn raikkaiden \VUgras{terälehtiensä} iloisen vastakohdan.

%%K t08-c1-03-1 p
3. --- Hedelmäpuut ovat koristautuneet lehdillä; hedelmät ovat kypsyneet. Pieni
\VUgras{kirsikkapuu} \textsuperscript{(4)} on täynnä kauniita
\VUgras{kirsikoita} \textsuperscript{(5)}, joita lapset suurella mielihyvällä
poimivat ja syövät.

%%K t08-c1-04-1 p
4. --- Nyt karja voi viettää koko päivän ulkoilmassa.

%%K t08-c1-04-2 p
\VUgras{Karitsat} \textsuperscript{(6)} ovat kasvaneet ja niistä on tullut
kauniita \VUgras{uuhia} \textsuperscript{(7)} ja kauniita \VUgras{pässejä}
\textsuperscript{(8)}, joilla on tuuhea villa. Ne syövät rauhallisesti
\VUgras{ruohoa} \VUgras{laitumella} \textsuperscript{(9)}, joka on täplikäs
\VUgras{kaunokaisista}.

%%K t08-c1-04-3 p
Pian aletaan keritä noita eläimiä saadakseen niiden \VUgras{villan}, josta
tehdään kangas vaatteisiimme.

%%K t08-tit-2 sub
\VUsaut{5.0mm}
\VUpk{10.2pt}{40}{Paimen.}
\VUsaut{3.4mm}

%%K t08-c1-05-1 p
5. --- \VUgras{Paimen} \textsuperscript{(10)} ei näytä pitävän niitä kovin
tarkasti silmällä. Hän välittää vain ukkosesta, joka kulkee ohi
taivaanrannalla, ja \VUgras{lammaspaimen} \textsuperscript{(13)}, joka vie
\VUgras{kenttäpullon} \textsuperscript{(14)} huulilleen, näyttää vieläkin
huolettomammalta.

%%K t08-c1-06-1 p
6. --- Helle on painostava; sillä \VUgras{koira} \textsuperscript{(15)} tulee,
hänkin, sammuttamaan janoaan; hän latkii \VUgras{puron} \textsuperscript{(16)}
raikasta vettä ja säikäyttää poloisen \VUgras{pienen sammakon}
\textsuperscript{(17)}, joka oli piiloutunut rantaruohikkoon. Tämä hyppää
kirkkaaseen veteen, jossa \VUgras{lumpeet} \textsuperscript{(18)} kukkivat.

%%K t08-c1-07-1 p
7. --- \VUgras{Västäräkki} \textsuperscript{(19)} kylpee pienen \VUgras{lähteen}
\textsuperscript{(20)} luona, joka pulppuaa kahden kallion välistä ja
piiloutuu \VUgras{orjantappurapensaiden} \textsuperscript{(21)} ja
\VUgras{orapihlajapensaiden} \textsuperscript{(22)} alle.

%%K t08-tit-3 sub
\VUsaut{5.0mm}
\VUpk{10.2pt}{40}{Elonkorjuu.}
\VUsaut{3.4mm}

%%K t08-c1-08-1 p
8. --- Maalaislapsille viljankorjuu on hyvin hauskaa aikaa. He tekevät iloisesti
\VUgras{kuperkeikkoja} \textsuperscript{(24)} \VUgras{olkikasoilla}
\textsuperscript{(23)}, he auttavat \VUgras{elonkorjaajia}
\textsuperscript{(26)}, ja samalla kun nämä niittävät \VUgras{viljapeltoa}
\textsuperscript{(27)} \VUgras{viikatteillaan} \textsuperscript{(25)}, jotka on
hyvin teroitettu \VUgras{kovasimella} \textsuperscript{(30)}, he kokoavat viljan
ja sitovat sen \VUgras{lyhteiksi} \textsuperscript{(31)} tai \VUgras{pieniksi
kimpuiksi} \textsuperscript{(32)}.

%%K t08-c1-09-1 p
9. --- Joskus heistä suurimpien annetaan leikata \VUgras{sirpillä}
\textsuperscript{(12)} \VUgras{kultaiset korret} \textsuperscript{(28)}, jotka
kantavat raskaita \VUgras{tähkiä} \textsuperscript{(29)} täynnä jyviä; ja pienet,
jotka vartioivat \VUgras{karjaa} \textsuperscript{(33)}, joka laiduntaa
\VUgras{laitumella} \textsuperscript{(34)} suuren \VUgras{lehmuksen}
\textsuperscript{(35)} varjossa, pyydystävät \VUgras{haavillaan}
\textsuperscript{(36)} kauniita \VUgras{perhosia}.

%%K t08-c1-09-2 p
Muutamat kiipeävät jopa \VUgras{kärryille} \textsuperscript{(37)}
järjestääkseen lyhteet, joita heille ojennetaan \VUgras{hangolla}
\textsuperscript{(38)}.

%%K t08-c1-10-1 p
10. --- Koko \VUgras{tasangolla} \textsuperscript{(39)}, jossa \VUgras{leivoset}
\textsuperscript{(41)} alkavat lentää ylös, vallitsee vilkkaus. Kaikki ovat
askaroimassa elonkorjuun parissa. Mutta samalla kun köyhät jatkavat vanhojen
työkalujen käyttöä, \VUgras{viikatteiden, sirppien} ja \VUgras{varstojen}
\textsuperscript{(40)}, rikkaat maanomistajat antavat niittää vehnänsä tai
\VUgras{ruispeltonsa} \textsuperscript{(43)} \VUgras{niittokoneella}
\textsuperscript{(42)}. Sitten tehdään, koska lyhteitä ei tarvitse kantaa
latoon, suuria \VUgras{lyhdekasoja} \textsuperscript{(44)}.

%%K t08-c1-11-1 p
11. --- Silloin sinne tuodaan \VUgras{puimakone} \textsuperscript{(47)}, jota
\VUgras{lokomobiili} \textsuperscript{(45)} käyttää \VUgras{käyttöhihnalla}
\textsuperscript{(46)}, ja niin jyvä erotetaan \VUgras{oljesta} lähtemättä siltä
pellolta, jolle se kylvettiin.

%%K t08-tit-4 sub
\VUsaut{5.0mm}
\VUpk{10.2pt}{40}{Ukonilma.}
\VUsaut{3.4mm}

%%K t08-c1-12-1 p
12. --- Usein tulee ankaria \VUgras{ukonilmoja} \textsuperscript{(53)}
elonkorjuun aikana. Ensin kuuluu \VUgras{ukkosen} jyrinä kaukaa; \VUgras{raju
länsituuli} \textsuperscript{(54)} alkaa puhaltaa ja taivuttaa puhkuen
\VUgras{metsän} \textsuperscript{(49)} paksuimpia puita. Mustat
\VUgras{pilvet} \textsuperscript{(56)} valtaavat taivaan; niitä halkovat
\VUgras{salaman} \textsuperscript{(55)} häikäisevät siksakit. \VUgras{Sade} ja
joskus \VUgras{rakeet} alkavat pudota ja lyövät maahan sadon, joka on valmis
leikattavaksi.

%%K t08-c1-13-1 p
13. --- Usein salama sytyttää rakennuksen. Niinpä viime vuonna
\VUgras{tuulimyllyn} \textsuperscript{(50)} katto pantiin tuhkaksi ja kaksi sen
\VUgras{siipeä} \textsuperscript{(51)} murskattiin.

%%K t08-c1-14-1 p
14. --- Muutaman tunnin kuluttua rauha palaa. Säteilevä \VUgras{sateenkaari}
\textsuperscript{(52)} ilmestyy taivaanrannalle ja luonto ottaa takaisin
elämänsä, joka keskeytyi muutamaksi hetkeksi. Maalaiset, jotka olivat paenneet
\VUgras{mökkeihin} \textsuperscript{(48)}, alkavat taas työskennellä ja koettavat
urhoollisesti korjata ne vahingot, jotka he ovat juuri kärsineet.

%%K t08-tit-5 sub
\VUsaut{6.0mm}
\VUcentre{13.2pt}{90}{\textit{Toinen kohtaus.}}
\VUsaut{3.6mm}

%%K t08-tit-6 sub
\VUpk{10.2pt}{40}{Metsästys. --- Viininkorjuu.}
\VUsaut{1.4mm}
\VUpk{10.2pt}{40}{Kalastus.}
\VUsaut{4.0mm}

%%K t08-c2-01-1 p
1. --- \VUgras{Elokuun} lopulla tai \VUgras{syyskuun} puolivälissä, seudun mukaan,
tapahtuu metsästyksen avaus. Tuskin ovat ensimmäiset auringonsäteet saaneet
\VUgras{sisiliskot} \textsuperscript{(2)} tulemaan ulos kolostaan, kun jo
\VUgras{metsästäjä} \textsuperscript{(5)} lähtee matkaan \VUgras{koirineen}
\textsuperscript{(4)}.

%%K t08-c2-02-1 p
2. --- Hän on pukenut ylleen tukevan \VUgras{metsästyspukunsa}
\textsuperscript{(9)}, joka on paksua kangasta, pannut jalkaansa nahkaiset
\VUgras{säärystimet}, joilla hän onnistuu kulkemaan kosteassa ruohikossa, jossa
\VUgras{rupikonnat} \textsuperscript{(1)} piileskelevät; hän on ottanut
\VUgras{kiväärinsä} \textsuperscript{(6)} ja \VUgras{metsästyslaukkunsa}
\textsuperscript{(10)}, ja tässä hän lähtee matkaan.

%%K t08-c2-03-1 p
3. --- Takaa-ajon into vie hänet keskelle viinitarhaa, jossa viininkorjuu on
hyvin vilkasta. Viininviljelijän lapset, jotka tavallisesti vartioivat vuohia
tien varrella, antavat niiden tänään laiduntaa umpimähkään, ja sidottuaan vanhan
\VUgras{pukin} \textsuperscript{(3)} paaluun, joka ei jättäisi käyttämättä
vapauttaan karatakseen, he ottavat vuorostaan osansa huvista. Toinen tarttuu
rypäleterttuun \VUgras{korjuukorissa} \textsuperscript{(12)} ja huvittelee
antamalla \VUgras{mehun} \textsuperscript{(14)} valua pieneen pikariin
\textsuperscript{(13)} tehdäkseen rypälemehua (käymätöntä viiniä). Toinen
koristautuu \VUgras{lehviseppeleellä} \textsuperscript{(15)}, jonka hän on juuri
punonut. Heidän vieressään röyhkeät \VUgras{varpuset} \textsuperscript{(16)}
tulevat aivan puristimen ääreen ryöstämään rypäleitä.

%%K t08-c2-04-1 p
4. --- Samalla kun lapset huvittelevat, miehet työskentelevät.
\VUgras{Viininkorjaajat} \textsuperscript{(18)} kantavat rypäleet kellariin
\VUgras{selkäkoreilla} \textsuperscript{(17)}; \VUgras{tynnyrintekijä}
\textsuperscript{(19)} korjaa \VUgras{tynnyreitä} \textsuperscript{(20)},
tarkistaa ovatko \VUgras{kimmet} \textsuperscript{(22)} ja \VUgras{pohjat}
\textsuperscript{(21)} hyvässä kunnossa, ja vaihtaa katkenneet \VUgras{vanteet}
\textsuperscript{(23)}. Hän on myös työskennellyt \VUgras{sammioiden}
\textsuperscript{(25)} parissa, joihin kaadetaan \VUgras{rypäleet}
\textsuperscript{(26)} niiden käymään panemiseksi.

%%K t08-c2-05-1 p
5. --- Kun käyminen on päättynyt, viini lasketaan ja jäännös pannaan
\VUgras{puristimeen} \textsuperscript{(28)}, ja kiristämällä vähitellen
\VUgras{suurta ruuvia} \textsuperscript{(29)} puristetaan ulos mehu, joka valuu
\VUgras{pieneen sammioon} \textsuperscript{(32)}, ja saadaan vielä
\VUgras{viiniä} \textsuperscript{(31)}.

%%K t08-tit-8 sub
\VUsaut{6.0mm}
\VUpk{10.2pt}{40}{Olut ja siideri.}
\VUsaut{3.4mm}

%%K t08-c2-06-1 p
6. --- \VUgras{Kellarin} \textsuperscript{(27)} muuria vasten kiipeää
\VUgras{humalan} \textsuperscript{(30)} oksa. Siellä se on vain koristekasvi,
mutta joissakin maissa, joissa viiniköynnös on hyvin harvinainen, humalaa
viljellään \VUgras{oluen} tekemiseksi.

%%K t08-c2-07-1 p
7. --- Pieni \VUgras{omenapuu} \textsuperscript{(33)} on lastattu omenilla,
jotka, kun ne ovat kypsiä, antavat erinomaista \VUgras{siideriä}.
\VUgras{Häkissään} \textsuperscript{(34)} \VUgras{kanarialintu}
\textsuperscript{(35)} ja \VUgras{tikli} \textsuperscript{(36)} katsovat
kateellisin silmin varpusia, jotka leikkivät vapaina.

%%K t08-c2-08-1 p
8. --- Niin kauas kuin silmä kantaa, kukkuloiden rinteellä on \VUgras{viinitappeja}
\textsuperscript{(40)} rivissä, jotka kannattavat komeita \VUgras{runkoja}
\textsuperscript{(39)}, raskaiden \VUgras{terttujen} koristamia.

%%K t08-c2-08-2 p
Koko vuoden \VUgras{viininviljelijä} \textsuperscript{(42)} on työskennellyt
hoitaakseen \VUgras{viinitarhaansa} \textsuperscript{(38)}, ja nyt hänet
palkitaan vaivoistaan kauniilla sadolla. \VUgras{Viininkorjaajattaret}
\textsuperscript{(41)} täyttävät sammiot, jotka ovat kärryillä, joita
\VUgras{muulit} \textsuperscript{(43)} vetävät, ja kaikki ovat iloisia.

%%K t08-tit-9 sub
\VUsaut{5.0mm}
\VUpk{10.2pt}{40}{Kalastus.}
\VUsaut{3.4mm}

%%K t08-c2-09-1 p
9. --- Niin on myös \VUgras{onkijoiden} \textsuperscript{(44)} laita, jotka
aamusta asti tulivat asettumaan veden rannalle \VUgras{onkivapoineen}
\textsuperscript{(45)}.

%%K t08-c2-09-2 p
\VUgras{Kalat} \textsuperscript{(47)} nappaavat \VUgras{koukkuun} heti kun he
kastavat \VUgras{siimansa} \textsuperscript{(46)} veteen, ja tuskin he ehtivät
uusia \VUgras{syöttiä}, kun jo uusi uhri sätkii siiman päässä.

%%K t08-c2-09-3 p
Kuitenkin \VUgras{verkkokalastaja} \textsuperscript{(48)}, joka
\VUgras{veneestään} \textsuperscript{(49)} heittää \VUgras{heittoverkon}
keskelle \VUgras{jokea} \textsuperscript{(50)}, pyydystää enemmän kaloja.

%%K t08-c2-10-1 p
10. --- Syksy on kauniiden kävelyretkien vuodenaika : jalan, \VUgras{hevosen}
\textsuperscript{(52)} selässä, \VUgras{polkupyörällä} \textsuperscript{(53)},
\VUgras{autolla} \textsuperscript{(54)}. \VUgras{Tiet} \textsuperscript{(51)}
ovat puhtaat, ja nautitaan viimeisistä kauniista päivistä, sillä tässä on, että
\VUgras{pääskyset} \textsuperscript{(56)} jo kokoontuvat katoille lentääkseen
kaikin siivin lämpimämpiin maihin. Vanhan \VUgras{saksanpähkinäpuun}
\textsuperscript{(57)} lehvistö kellastuu, \VUgras{pähkinät} putoavat ja korkeat
\VUgras{puut}, jotka ovat ryhmittyneinä \VUgras{lehtona}
\textsuperscript{(58)} \VUgras{kukkulalla} \textsuperscript{(59)}, pudottavat
pian lehtensä.

%%K t08-c2-11-1 p
11. --- \VUgras{Aurinko} \textsuperscript{(60)} nousee myöhemmin, päivät
lyhenevät, jokseenkin purevaa kylmyyttä alkaa tuntua \VUgras{aamulla}, ja tässä
on, että jo parvet \VUgras{lehtokurppia} \textsuperscript{(61)} jättävät
epävieraanvaraiset seutumme.
\VUsaut{6.0mm}
\VUfilet{22mm}
